\documentclass[a4paper]{article}
\usepackage{graphicx}
\usepackage{amsmath,amssymb}
\usepackage{minted}
\usepackage{multirow}
\usepackage{float}
\usepackage{algorithm}
\usepackage{algpseudocode}
\usepackage{todonotes}
\usepackage{forest}
\usepackage{xstring}
\usepackage{xcolor}
\usepackage[most]{tcolorbox}
\usepackage{xparse}
\usepackage{natbib}
\usetikzlibrary{graphs,graphdrawing}
\usegdlibrary{layered}
% ND Mind Mirror default input files.
\input{/home/donkarlo/Dropbox/repo/nd_language_ai_project/src/nd_language_ai/formal/computer/programming/tex/latex/code_clips/external_graphviz.tex}
\input{/home/donkarlo/Dropbox/repo/nd_language_ai_project/src/nd_language_ai/formal/computer/programming/tex/latex/part/control_sequence/kind/semtag/semtag.tex}
\input{/home/donkarlo/Dropbox/repo/nd_language_ai_project/src/nd_language_ai/formal/computer/programming/tex/latex/code_clips/seehere.tex}

\title{Linear Regression for Trend Detection}
\date{}
\author{Mohammad Rahmani}
\begin{document}
\maketitle

\section{Definition}
Linear regression models a series as a linear function of elapsed time. The slope estimates the average change in a quantity per unit of time.

\section{Model}
For observations $(t,y_t)$, use
\[
 y_t=\beta_0+\beta_1 t+\varepsilon_t.
\]
The ordinary least-squares (OLS) estimate of the slope is
\[
 \widehat\beta_1 =
 \frac{\sum_t(t-\overline t)(y_t-\overline y)}
      {\sum_t(t-\overline t)^2}.
\]
A positive estimated slope indicates fitted growth, a negative slope indicates fitted decline. The slope sign alone does \emph{not} establish a statistically significant effect.

\section{Example: E-Commerce Revenue}
A fitted slope of 45 euros/day corresponds to an average model-based increase of 45 euros of revenue for each extra day. It is not proof that the passage of time caused the increase.

\section{Accounting for Other Drivers}
A retail trend model may require day-of-week indicators $D_{j,t}$ and a promotion indicator $P_t$:
\[
 y_t=\beta_0+\beta_1t+\sum_{j=1}^{6}\gamma_jD_{j,t}
       +\delta P_t+\varepsilon_t.
\]
Interpret $\beta_1$ conditional on these covariates. Other potentially important factors include product mix, price changes, stockouts, and holidays.

\section{Statistical Inference}
A conventional hypothesis is $H_0:\beta_1=0$. Standard OLS tests require an appropriate error structure. Serial correlation or changing residual variance can invalidate naive standard errors. When needed, use valid time-series error modeling or robust standard errors.

\section{Advantages}
\begin{itemize}
 \item Interpretable slope with practical units.
 \item Simple baseline, easily extended with explanatory variables.
 \item Confidence intervals possible under appropriate assumptions.
\end{itemize}
\section{Limitations}
\begin{itemize}
 \item One global slope can hide nonlinear changes and breakpoints.
 \item OLS is sensitive to outliers.
 \item Forecast extrapolation beyond the observed period can be unreliable.
\end{itemize}
\section{Practical Guidance}
Inspect residual autocorrelation and the scatterplot. Compare with Theil--Sen for robustness and use LOESS or segmented models if the trend changes shape.

\bibliographystyle{plainnat}
\bibliography{/home/donkarlo/Dropbox/repo/nd_shared_research/bibliography/bibliography.bib}
\end{document}
